\subsection{Limitations}\label{limitations}

The compatibility premise is assumed, not derived. The central bridge is
a recording-level compatibility condition between the latent assignment
and the population K-means objective, and the closed level-set and
quadratic-growth conditions are hypotheses as well. Nothing in the
generative model forces them to hold, and the earlier exchange argument
that sought to derive them is withdrawn. The hard-ball and separation
assumptions are an idealized benchmark, not a description of the EEG
data, and the paper makes no claim that the covariance features satisfy
them. Proving alignment on a structured class of distributions is the
natural next step, and the quantitative growth bound is local, with no
global version claimed.

The number of environments is not identified. Every theoretical
statement conditions on \(K\), and there is no proof that \(K\) is
identifiable or consistently estimable; on the pooled real data the
blind selector's choice is unstable, alternating between two and three
clusters across seeds, while the forced two-cluster merge is the stable
pattern. Automatic selection, full-pool standardization, and
finite-initialization Lloyd stay outside the unconditional guarantees,
and the window-level and recording-level estimands coincide only under
extra conditions.

Several technical pieces remain open. An algorithm-independent bound on
the training-side center error is not established, and the projection
cross-terms are treated through an all-direction control, not a closed
form. The exact correspondence between estimated and population objects,
the unification of the scaler coordinates, and the fitted center term
are each stated as conditions, and the K-means basin of attraction and a
local-to-global proof are not provided. On the data side, bad-channel
interpolation and an absolute quality threshold are absent, and the
channel-sensitivity result rules out only a single-channel explanation,
not distributed artifacts or a shared device fingerprint.

The empirical and identification claims are bounded. The real-data label
is a technical mixture of device, task, population, duration, reference,
and preprocessing whose driver cannot be separated, and no biological or
causal reading is available. The three-source design gives repeatability
evidence, not a generality proof, and full cross-domain generalization
is not established. The partial-identification framing is descriptive,
not a proved sharp set, the resampling intervals carry no coverage
guarantee, and the general non-Gaussian, fully blind problem remains
open.

\subsection{Conclusion}\label{conclusion}

The paper asked whether a covariance-based environment structure stays
identifiable once the labels disappear and the observations are pooled.
The central claim is a single one. Under unlabeled pooling,
covariance-based environment discovery has a sharp identifiability
boundary and a conditional recovery boundary, so source-associated
technical structure is recoverable at a fixed number of environments
while blind selection of that number and any causal reading remain
unreliable. The results that follow are the components of this claim,
not a set of independent answers.

A two-dimensional witness shows that two genuinely different partitions
can induce the same pooled distribution, enough to rule out treating
identifiability as automatic without making a universal claim. In a
constrained, whitened two-component frame the component covariance pair
is point-identified, plainly a special case, while the moment-order
threshold beyond it is reported as cited work. Relax the assumptions and
the target can only be described as a set of consistent structures, not
a proved sharp set.

The practical contribution is a recording-level framework that keeps the
observation mechanism, the estimand, and the algorithm version distinct.
Under fixed \(K\) and a pre-given transformation, the global empirical
K-means solution carries a finite-sample risk-gap guarantee, and a
sample-split variant gives a held-out decomposition whose training-side
term stays explicit. The standard components are not presented as novel.
On three independent EEG sources, source-associated structure is
recoverable under a known \(K\) and a balanced design, and the
subject-cluster bootstrap stays centered above zero as a diagnostic
without a coverage guarantee. Three boundary findings accompany this:
the blind selector's cluster-number choice is unstable, the unbalanced
plug-in is weak, and within-source recovery is absent.

The two parts answer complementary questions. Pooled-distribution
identifiability asks whether the population law determines the latent
decomposition at all, and recording-level recovery asks whether a chosen
clustering objective can approximate the latent labels under additional
assumptions. Identifiability of the pooled model does not by itself
imply recording-level clustering recovery, and a successful clustering
analysis does not establish causal-environment or causal-representation
identification.

The recovered recording-level groups should be interpreted as candidate
environments rather than validated causal environments. We do not
establish that these groups identify a causal representation, an
invariant mechanism, or a causal graph. A natural next step is to
propagate uncertainty in the recovered grouping into downstream causal
representation learning and to test whether the resulting groups improve
cross-environment generalization or causal identifiability.

Deriving the compatibility condition from a structured distribution
class, establishing the identifiability of \(K\), closing the
training-side bound, and extending the framework to non-Gaussian and
unseen settings would each move the boundary outward. Until then, the
claims are offered with their conditions attached. The complete
implementation and the scripts that reproduce every table and figure are
released in a public repository, and the data sources and access details
are stated in the Data Availability section.
